% !TeX program = xelatex
%
% Original starter for latextemplates.com -- not upstream's own main.tex,
% which is a tutorial deck documenting the theme's own macros (and uses
% minted, which needs -shell-escape + pygments, unavailable in this
% compile pipeline). Only the theme files and background/logo images are
% vendored from github.com/Thysrael/GreenParkBeamerTheme.
\documentclass{ctexbeamer}

\usetheme{gpark}
\usepackage{booktabs}

\title{无人机集群协同路径规划}
\subtitle{北航六系 2026 年秋季组会报告}
\author{李明轩}
\date{\today}
\titlebackground{images/titlebackground}
\background{images/background2}

\begin{document}

\maketitle

\begin{frame}{目录}
  \tableofcontents
\end{frame}

\section{研究背景}

\begin{frame}{问题描述}
  \begin{itemize}
    \item 多无人机协同任务对实时路径规划提出更高要求
    \item 传统集中式规划在节点数增加时计算开销急剧上升
    \item 本报告提出一种分布式局部规划方法
  \end{itemize}
\end{frame}

\begin{frame}{相关工作对比}
  \begin{table}
    \centering
    \begin{tabular}{lrr}
      \toprule
      方法 & 节点规模 & 平均耗时 (ms) \\
      \midrule
      集中式 A* & 50 & 820 \\
      分布式 RRT & 50 & 340 \\
      本文方法 & 50 & \textbf{126} \\
      \bottomrule
    \end{tabular}
    \caption{三种方法在相同场景下的耗时对比}
  \end{table}
\end{frame}

\section{方法}

\begin{frame}{总体框架}
  \begin{itemize}
    \item 每个无人机仅依赖局部通信半径内的邻居信息
    \item 采用势场法结合局部避碰约束生成候选路径
    \item 通过一致性协议在邻居间同步路径意图，避免死锁
  \end{itemize}
\end{frame}

\begin{frame}{关键公式}
  路径代价函数定义为：
  \[
    J(\mathbf{p}) = \alpha \cdot d(\mathbf{p}, \mathbf{p}_{\text{goal}}) + \beta \cdot \sum_{i \in \mathcal{N}} \phi(\mathbf{p}, \mathbf{p}_i)
  \]
  其中 $\phi$ 为邻居间的排斥势场，$\alpha, \beta$ 为可调权重。
\end{frame}

\section{实验结果}

\begin{frame}{仿真设置}
  \begin{itemize}
    \item 仿真环境：50 架无人机，随机分布起点与目标点
    \item 通信半径：15 米，更新频率：10 Hz
    \item 对比集中式 A* 与分布式 RRT 两种基线方法
  \end{itemize}
\end{frame}

\begin{frame}{成功率与耗时}
  \begin{table}
    \centering
    \begin{tabular}{lrr}
      \toprule
      节点规模 & 成功率 & 平均耗时 (ms) \\
      \midrule
      20 & 99.2\% & 58 \\
      50 & 97.6\% & 126 \\
      100 & 94.1\% & 241 \\
      \bottomrule
    \end{tabular}
    \caption{不同规模下的成功率与耗时}
  \end{table}
\end{frame}

\section{总结}

\begin{frame}{结论与展望}
  \begin{itemize}
    \item 分布式方法在保持较高成功率的同时显著降低了耗时
    \item 通信半径是影响成功率的关键参数，仍需进一步调优
    \item 后续工作将扩展至动态障碍物场景
  \end{itemize}
\end{frame}

\begin{frame}{谢谢聆听}
  \centering
  {\Large 谢谢！欢迎提问}\\[1em]
  li.mingxuan@buaa.edu.cn
\end{frame}

\end{document}
